% G3-1 Alur penelitian.
% Sumber: diolah penulis; rincian collector dari
% context/konteks_collector_data_v3.0.md.
\begin{tikzpicture}[x=1mm,y=-1mm,font=\footnotesize,
    tahap/.style={draw,rectangle,text width=96mm,align=left,
      inner sep=3pt,minimum height=17mm},
    nomor/.style={draw,circle,inner sep=0pt,minimum size=7mm,font=\small},
    panah/.style={-{Stealth[length=2.2mm]}}]

  % Tinggi kotak 17mm dan jarak antarbaris 23mm, seragam untuk semua langkah.
  \node[nomor] at (6,0) {1};
  \node[tahap] (t1) at (62,0)
    {\textbf{Pengumpulan data.} Candle 1 menit BTC/IDR dari API publik
     Indodax, diambil otomatis setiap menit oleh \textit{collector} di VPS
     dan disimpan di basis data SQLite.};

  \node[nomor] at (6,23) {2};
  \node[tahap] (t2) at (62,23)
    {\textbf{Praproses.} Candle 5 menit dibentuk dari agregasi candle
     1 menit. Harga \textit{typical}
     $(\text{high}+\text{low}+\text{close})/3$ dipakai untuk kedua dataset.
     Bar yang hilang dibiarkan kosong, tidak diisi.};

  \node[nomor] at (6,46) {3};
  \node[tahap] (t3) at (62,46)
    {\textbf{Pembentukan fitur dan target.} Fitur hanya dari data sampai
     waktu $t$; target $r_{t,h} = \ln(P_{t+h}/P_t)$.};

  \node[nomor] at (6,69) {4};
  \node[tahap] (t4) at (62,69)
    {\textbf{Pembagian data} berurutan waktu, dengan \textit{purging}
     $h$ bar terakhir pada data latih.};

  \node[nomor] at (6,92) {5};
  \node[tahap] (t5) at (62,92)
    {\textbf{Optimasi \textit{hyperparameter}} dengan validasi
     \textit{walk-forward} di dalam data latih.};

  \node[nomor] at (6,115) {6};
  \node[tahap] (t6) at (62,115)
    {\textbf{Pelatihan model akhir dan prediksi.} Return yang diramal
     dikonversi ke rupiah: $\hat{P}_{t+h} = P_t\exp(\hat{r}_{t,h})$.};

  \node[nomor] at (6,138) {7};
  \node[tahap] (t7) at (62,138)
    {\textbf{Evaluasi terhadap Naive.} Metrik MAE, RMSE, MAPE, akurasi arah,
     dan uji Diebold-Mariano dengan koreksi HLN.};

  \foreach \a/\b in {1/2,2/3,3/4,4/5,5/6,6/7}{
    \draw[panah] (t\a.south) -- (t\b.north);
  }
\end{tikzpicture}
